% Standalone problem document, generated from the adjacent README.md.
% Compile with XeLaTeX. Mathematical content is maintained in Markdown.
\documentclass[11pt,a4paper]{article}
\usepackage[fontsize=10.5pt]{scrextend}
\usepackage[margin=22mm,headheight=20pt,headsep=9mm,footskip=11mm]{geometry}
\usepackage{fontspec}
\setmainfont{texgyrepagella}[Extension=.otf,UprightFont=*-regular,BoldFont=*-bold,ItalicFont=*-italic,BoldItalicFont=*-bolditalic]
\setsansfont{texgyreheros}[Extension=.otf,UprightFont=*-regular,BoldFont=*-bold,ItalicFont=*-italic,BoldItalicFont=*-bolditalic]
\setmonofont{texgyrecursor}[Extension=.otf,UprightFont=*-regular,BoldFont=*-bold,ItalicFont=*-italic,BoldItalicFont=*-bolditalic,Scale=MatchLowercase]
\usepackage{amsmath,amssymb}
\usepackage{unicode-math}
\setmathfont{texgyrepagella-math.otf}
\usepackage{xcolor,microtype,enumitem,needspace,fancyhdr}
\usepackage{bookmark}
\definecolor{ink}{HTML}{17344B}
\definecolor{accent}{HTML}{197D80}
\definecolor{muted}{HTML}{596773}
\hypersetup{colorlinks=true,linkcolor=accent,urlcolor=accent,pdftitle={TR-27 - Border-rank
deficiency forcing strict submultiplicativity at the tensor
square},pdfauthor={Open Problems in Numerical Linear Algebra}}
\urlstyle{same}
\setlength{\parindent}{0pt}
\setlength{\parskip}{5pt plus 1pt minus 1pt}
\linespread{1.04}
\setlength{\emergencystretch}{3em}
\widowpenalty=10000
\clubpenalty=10000
\displaywidowpenalty=10000
\allowdisplaybreaks[1]
\setlist{nosep,leftmargin=1.5em}
\providecommand{\tightlist}{\setlength{\itemsep}{0pt}\setlength{\parskip}{0pt}}
\providecommand{\pandocbounded}[1]{#1}
\pagestyle{fancy}
\fancyhf{}
\fancyhead[L]{\sffamily\footnotesize\color{muted}OPEN PROBLEMS IN NUMERICAL LINEAR ALGEBRA}
\fancyhead[R]{\sffamily\bfseries\footnotesize\color{accent}TR-27}
\fancyfoot[L]{\parbox[b]{0.9\textwidth}{\sffamily\fontsize{8}{10}\selectfont\color{muted}Editorial ratings; a bounded search cannot certify that a problem remains unsolved.\\Verification check: 2026-09-22}}
\fancyfoot[R]{\sffamily\footnotesize\color{muted}\thepage}
\renewcommand{\headrulewidth}{0.3pt}
\renewcommand{\headrule}{\hbox to\headwidth{\color{accent}\leaders\hrule height \headrulewidth\hfill}}
\makeatletter
\renewcommand{\section}{\Needspace{5\baselineskip}\@startsection{section}{1}{0pt}{12pt plus 2pt minus 2pt}{4pt}{\sffamily\bfseries\large\color{ink}}}
\renewcommand{\subsection}{\Needspace{4\baselineskip}\@startsection{subsection}{2}{0pt}{9pt plus 1pt minus 1pt}{3pt}{\sffamily\bfseries\normalsize\color{ink}}}
\makeatother
\setcounter{secnumdepth}{0}
\begin{document}
{\sffamily\small\color{muted}Tensor computations\par}
\vspace{3pt}
{\raggedright\hyphenpenalty=10000\exhyphenpenalty=10000\sffamily\bfseries\fontsize{21}{24}\selectfont\color{ink}Border-rank
deficiency forcing strict submultiplicativity at the tensor square\par}
\vspace{7pt}
{\sffamily\small\textbf{Difficulty:} extreme\quad\textbf{Importance:} interesting
to the community\par}
{\sffamily\small\bfseries\color{accent}Status: Lean verified\par}
\vspace{4pt}
{\color{accent}\hrule height 0.8pt}
\vspace{6pt}
\textbf{Rating rationale:} Forcing a saving at exactly two copies for
every projective variety goes beyond the known eventual-power mechanism
and presents a general structural barrier. The relationship between
degeneration and repeated decomposition matters to the tensor-complexity
community.

\textbf{Negative resolution recorded 2026-09-11.} Matthew J. Colbrook's
\href{https://github.com/ajt60gaibb/OpenProblemsInNLA/blob/main/references/colbrook-tensor-metrics-rank-2026-09-11/manuscripts/tr27_solution.pdf}{complete
manuscript, Theorem 1 and Section 4}
(\href{https://github.com/ajt60gaibb/OpenProblemsInNLA/blob/main/references/colbrook-tensor-metrics-rank-2026-09-11/manuscripts/tr27_solution.tex}{LaTeX
source}) constructs a smooth, irreducible, reduced, nondegenerate
complex projective curve \(X\subset\mathbb P^{11}\) and a point \(p\)
with \(\underline R_X(p)=2\), \(R_X(p)=3\), and
\(R_{X\times X}(p\otimes p)=9\). This refutes the exact universal
implication below, using the specified Segre product with unrestricted
complex coefficients. More generally, the construction delays strict
submultiplicativity for any prescribed finite number of powers; the
variety depends on that prescribed delay. It does not refute versions
restricted to Segre or Veronese varieties or the known eventual-power
saving. The full original source passed an
\href{https://github.com/ajt60gaibb/OpenProblemsInNLA/blob/main/references/colbrook-tensor-metrics-rank-2026-09-11/verification/reviews/TR-27-review.md}{independent
mathematical agent review}.
\href{https://github.com/ajt60gaibb/OpenProblemsInNLA/blob/main/references/colbrook-tensor-metrics-rank-2026-09-11/README.md}{Submission
and verification record}. That 2026-09-11 review is documented
AI-assisted verification, not external human peer review; no priority
claim is made. The original target and former difficulty and importance
ratings are retained below as historical context.

\textbf{Author feedback --- 2026-09-17.} In personal correspondence with
Alex Townsend, Alessandra Bernardi, a coauthor of Conjecture 1.1, said
she thought the argument works and agreed that it contradicts the
conjecture in its stated generality. \textbf{The restrictions to Segre
and Veronese varieties remain open.} This is informal mathematical
feedback, not a formal referee report.
\href{https://github.com/ajt60gaibb/OpenProblemsInNLA/blob/main/references/colbrook-tensor-metrics-rank-2026-09-11/README.md\#tr-27-author-feedback--2026-09-17}{Correspondence
summary}.

\newpage

\subsection{Lean proof and verification evidence -
2026-09-22}\label{lean-proof-and-verification-evidence---2026-09-22}

\textbf{The complete original conjecture is Lean verified, with a
negative answer.} The
\href{https://github.com/ajt60gaibb/OpenProblemsInNLA/blob/775e8b169119c4045b07db7666eda8c001ae3bd1/tensor-computations/TR-27/lean/NLA/TR27/Counterexample.lean}{formal
proof at revision 775e8b1} constructs a reduced irreducible
nondegenerate complex projective curve and a point with rank 3, border
rank at most 2, and ordinary Segre tensor-square rank 9. These
inequalities refute the full implication below. Decompositions allow
arbitrary complex coefficients, repeated points, zero summands, the
point at infinity, and independently chosen left and right tensor
factors. The proof covers the whole variety, using an algebraic proof
that the parameter cone equals the entire closed zero locus.

\textbf{Lean formalization:} George Stepaniants, Department of Computing
and Mathematical Sciences, California Institute of Technology, with
substantial AI-agent assistance. \textbf{Matthew J. Colbrook} retains
attribution for the original counterexample and mathematical proof. The
conjecture authorship, manuscript, original statement and author
feedback above remain preserved.

The
\href{https://github.com/ajt60gaibb/OpenProblemsInNLA/blob/775e8b169119c4045b07db7666eda8c001ae3bd1/tensor-computations/TR-27/lean/Solution.lean}{25
checked declarations} establish the projective, Zariski-closure and
Segre semantics, the whole-image geometry, the rank and border-rank
bounds, and the unconditional final theorems
\textit{NLA.TR27.projective\_counterexample} and
\textit{NLA.TR27.original\_conjecture\_false}. The
\href{https://github.com/ajt60gaibb/OpenProblemsInNLA/blob/main/tensor-computations/TR-27/lean/Challenge.lean}{statement
boundary} was independently reviewed before implementation and remains
byte-identical to its frozen receipt. Two fresh nonauthor agents
approved the complete mathematical source and separately checked the
actual verification evidence:
\href{https://github.com/ajt60gaibb/OpenProblemsInNLA/blob/main/tensor-computations/TR-27/lean/reviews/final-fidelity-referee.md}{fidelity
source review} and
\href{https://github.com/ajt60gaibb/OpenProblemsInNLA/blob/main/tensor-computations/TR-27/lean/reviews/final-fidelity-completion.md}{completion
review}, and
\href{https://github.com/ajt60gaibb/OpenProblemsInNLA/blob/main/tensor-computations/TR-27/lean/reviews/final-correctness-referee.md}{correctness
source review} and
\href{https://github.com/ajt60gaibb/OpenProblemsInNLA/blob/main/tensor-computations/TR-27/lean/reviews/final-correctness-completion.md}{completion
review}. These are independent AI reviews adapting the repository's Tau
Ceti protocol, not external human peer review or official Tau Ceti
endorsement.

The
\href{https://github.com/ajt60gaibb/OpenProblemsInNLA/blob/main/tensor-computations/TR-27/lean/verification/linux/OPERATIONAL-REVIEW.md}{successful
non-root Linux run and operational audit} bind the immutable proof
inputs to all 25 successful Comparator comparisons, default-kernel
replay, authentic LeanCert kernel assertions, dependency revisions and
isolation/rejection controls. Every selected declaration uses only
\textit{propext}, \textit{Classical.choice} and \textit{Quot.sound}. The
pins are Lean 4.33.1, LeanCert
\textit{621a43d7cf21f87872392a01e874f2f1dbddc926} and Mathlib
\textit{0df444a360eaa60ab8c11dca51a86af692955474}; the
\href{https://github.com/ajt60gaibb/OpenProblemsInNLA/blob/main/tools/lean/source-lock.json}{shared
source lock} pins the Comparator and control tools. This exact algebraic
proof needs no numerical interval certificate.

The manuscript's additional smoothness, exact border rank two and
arbitrary prescribed finite-delay results retain their manuscript and
informal-review scope. They are unnecessary for this complete negative
answer and are not claimed as Lean verified. The Segre- and
Veronese-restricted questions remain open as recorded above. See the
\href{https://github.com/ajt60gaibb/OpenProblemsInNLA/blob/main/tensor-computations/TR-27/lean/README.md}{project
guide} and
\href{https://github.com/ajt60gaibb/OpenProblemsInNLA/blob/main/tensor-computations/TR-27/lean/formalization.yaml}{formalization
metadata}. On the documented
\href{https://github.com/ajt60gaibb/OpenProblemsInNLA/blob/main/docs/lean/README.md}{non-root
Linux environment}, reproduce from the repository root with:

\begin{verbatim}
tools/lean/bootstrap.sh /absolute/path/to/verification-tools
tools/lean/verify.sh \
  tensor-computations/TR-27/lean \
  /absolute/path/to/verification-tools
\end{verbatim}

\newpage

\subsection{Problem statement}\label{problem-statement}

Let \(V\) be a finite-dimensional complex vector space and let
\(X\subset\mathbb P(V)\) be an irreducible reduced nondegenerate complex
projective variety. Nondegenerate means that \(X\) spans
\(\mathbb P(V)\). For \(p\in\mathbb P(V)\), define \(R_X(p)\) as the
least number of points of \(X\) whose projective linear span contains
\(p\). Define \(\underline R_X(p)\) as the least \(s\) for which \(p\)
lies in the Zariski closure of \(\{q:R_X(q)\leq s\}\).

Embed \(X\times X\) in \(\mathbb P(V\otimes V)\) by the Segre map
\(([x],[y])\mapsto[x\otimes y]\). Is it always true that

\[\underline R_X(p)< R_X(p)
\quad\Longrightarrow\quad
R_{X\times X}(p\otimes p)< R_X(p)^2?\]

Here \(p\otimes p\) denotes the projective point represented by the
tensor square of any representative of \(p\). The product variety is
precisely the Segre image just defined. For tensor-rank examples, this
is the ordinary tensor product with both sets of modes retained; it does
not silently merge corresponding modes into a Kronecker product.

\subsection{Why it matters}\label{why-it-matters}

An exact tensor decomposition can require more terms than arbitrarily
close decompositions. The conjecture asks whether that deficiency
already yields a strict saving when two copies are decomposed together.
It connects ill-behaved rank limits with the cost of repeated
multilinear computation.

\subsection{References and status
check}\label{references-and-status-check}

\begin{enumerate}
\def\labelenumi{\arabic{enumi}.}
\tightlist
\item
  E. Ballico, A. Bernardi, F. Gesmundo, A. Oneto, and E. Ventura,
  \emph{Geometric conditions for strict submultiplicativity of rank and
  border rank},
  \href{https://arxiv.org/html/1909.03811v2}{arXiv:1909.03811v2},
  Conjecture 1.1; Annali di Matematica Pura ed Applicata \textbf{200}
  (2021), 187--210.
\item
  A. Oneto and E. Ventura, \emph{Ranks of tensors: geometry and
  applications},
  \href{https://doi.org/10.1007/s40574-025-00472-9}{Bollettino
  dell'Unione Matematica Italiana \textbf{18} (2025), 751--778},
  Conjecture 4.9.
\end{enumerate}

\subsubsection{Status check ---
2026-09-10}\label{status-check-2026-09-10}

Checked the original v2, Conjecture 1.1 and Theorem 4.1, the 2025
restatement, and targeted strict-submultiplicativity searches through
2026. Theorem 4.1 proves the implication for all binary forms, a
substantive subclass of the displayed target. The source also treats
ternary cubics in Proposition 4.2 and its proof. A saving at some
sufficiently large tensor power does not establish the required saving
at the square. No general proof or counterexample was located; the
product here retains both sets of factors.
\end{document}
